\documentclass[11pt]{article}

% Encoding & fonts (safer defaults)
\usepackage[utf8]{inputenc}
\usepackage[T1]{fontenc}
\usepackage{lmodern}

% Packages for formatting
\usepackage[a4paper,margin=1in]{geometry}
\usepackage{enumitem}
\usepackage{hyperref}
\usepackage{titlesec}
\usepackage{array}
\usepackage{xcolor}
\usepackage{booktabs} % nicer tables
\usepackage{amsmath}  % for math symbols like \rightarrow

% Section formatting
\titleformat{\section}{\large\bfseries}{}{0em}{}
\titleformat{\subsection}{\normalsize\bfseries}{}{0em}{}

% Header info
\newcommand{\university}{University of Vienna}
\newcommand{\coursename}{Applied Economics}
\newcommand{\coursecode}{040178-1 and 040178-2}
\newcommand{\semester}{Winter Term 2025/2026}
\newcommand{\instructor}{Pello Aspuru}
\newcommand{\email}{pello.aspuru@univie.ac.at}
\newcommand{\officehours}{By appointment, at pello.aspuru@univie.ac.at}
\newcommand{\office}{Department of Economics, Oskar-Morgenstern-Platz 1, 1090 Vienna, 05.610}
\newcommand{\courselanguage}{English}
\newcommand{\countECTS}{4 ECTS}

\begin{document}

% Title
\begin{center}
    {\Large \textbf{\coursename}} \\
    \coursecode \\
    \semester \\
    Language: \courselanguage \\
    \countECTS \\
    \vspace{0.3cm}
    {\university} \\
    \vspace{0.5cm}
    Instructor: \instructor \\
    Email: \email \\
    Office Hours: \officehours \\
    Office: \office \\
\end{center}

\section{Dates and Location}

Course code 040178-1
\begin{itemize}
    \item Wednesdays: 09:45--11:15, starting October 1, 2025.
    \item PC-Seminarraum 5, Oskar-Morgenstern-Platz 1, Floor: KG01, Room: U1.611
\end{itemize}

Course code 040178-2
\begin{itemize}
    \item Thursdays: 16:45--18:15, starting October 2, 2025.
    \item PC-Seminarraum 5, Oskar-Morgenstern-Platz 1, Floor: KG01, Room: U1.611
\end{itemize}

\section{Course Description}
This is a course in which students will reinforce the tools learned in Introductory Econometrics (BA) by applying them to real-world data sets using econometric software. Prior knowledge at that level will be assumed throughout. The aim is for students to get hands-on experience analyzing observational data using the open-source programming language R and to acquire the tools to carry out empirical research projects independently.

\section{Course Objectives}
In this course you will learn to:
\begin{itemize}
\item Use econometric tools to produce empirical analyses.
\item Interpret and critically evaluate existing empirical studies.
\item Understand possibilities and limitations of empirical analyses.
\end{itemize}

\section{Required Software}
\begin{itemize}
    \item Install R from \url{https://cran.r-project.org/}
    \item  Install RStudio from \url{https://posit.co/download/rstudio-desktop/}
\end{itemize}

\section{Course Schedule} This is a tentative and non-binding course schedule: 
\begin{itemize} 
    \item Week 1 (Oct 1st / 2nd): Introduction 
    \item Week 2 (Oct 8th / 9th): Multiple Linear Regression 1/2 
    \item[\textcolor{red}{$\rightarrow$}] \textcolor{red}{Deregistration period ends} 
    \item Week 3 (Oct 15th / 16th): Multiple Linear Regression 2/2 
    \item Week 4 (Oct 22nd / 23rd): Ordinary Least Squares (OLS) 
    \item Week 5 (Oct 29th / 30th): Inference 1/2 
    \item Week 6 (Nov 5th / 6th): Inference 2/2 
    \item[\textcolor{red}{$\rightarrow$}] \textcolor{red}{Week 7 (November 12th / 13th): Midterm exam} 
    \item Week 8 (Nov 19th / 20th): Causality and Correlation 
    \item Week 9 (Nov 26th / 27th): Instrumental Variables (IV) 
    \item Week 10 (Dec 3rd / 4th): Panel Data 
    \item[\textcolor{red}{$\rightarrow$}] \textcolor{red}{Dec 5th, 23:59: Problem set online.}
    \item Week 11 (Dec 10th / 11th): Differences in Differences (DD) 
    \item Week 12 (Dec 17th / 18th): Regression Discountinuity Design (RDD) 
    \item[\textcolor{red}{$\rightarrow$}] \textcolor{red}{Dec 19th, 23:59: Problem set submission deadline.}
    \item Week 13 (Jan 7th / 8th): TBD 
    \item Week 14 (Jan 14th / 15th): TBD 
    \item Week 15 (Jan 21st / 22nd): TBD 
    \item[\textcolor{red}{$\rightarrow$}] \textcolor{red}{Jan 29th, 16:45-18:15: Final exam}
\end{itemize}

\section{Grading Policy} %\begin{center} 
\begin{tabular}{lc} 
    Class Participation \quad \quad & 10\% \\ 
    Problem Set & 25\% \\ 
    Midterm Exam & 25\% \\ 
    Final Exam & 40\% \\ 
\end{tabular} 
    %\end{center}

\section{Problem Set}
\begin{itemize}
    \item The problem set can be submitted in teams of up to three students.
    \item To register your problem set team:
    \begin{enumerate}
        \item Enter \href{https://docs.google.com/spreadsheets/d/1EELY5wWKEA_pHf_HXenjt4S6EhgivH2ieC7nghsS0EE/edit?usp=sharing}{this} Google Sheet.
        \item Select your weekday tab.
        \item Add the team member names.  
    \end{enumerate}
\end{itemize}
% Add brief about scope, submission format (PDF via LMS), late policy, collaboration rules, etc.

\section{Academic Integrity}
Students are expected to uphold the highest standards of academic honesty. Plagiarism or cheating will result in disciplinary action in accordance with university policy.

\section{Accessibility}
Students requiring accommodations should contact the instructor and the university’s accessibility services.

\section{Useful References}
  \begin{enumerate}
    \item \textbf{Oswald} F, Viers V, Villedieu P, Kennedy G (2020).
Introduction to Econometrics with
R. SciencesPo Department of
Economics, Paris, France. \\
available online: \href{https://scpoecon.github.io/ScPoEconometrics/}{https://scpoecon.github.io/ScPoEconometrics/}.
   \item \textbf{Heiss, F. (2020).} Using R for Introductory Econometrics (2nd edition) 

    available online: \href{http://urfie.net/readOnline.html}{http://urfie.net/readOnline.html}
    \item \textbf{Hanck}, C., Arnold, M., Gerber A., and Schmelzer M. (2023). Introduction to Econometrics with R 

        available online: \href{https://www.econometrics-with-r.org/index.html}{https://www.econometrics-with-r.org/index.html}
  \end{enumerate}

\end{document}
